\documentclass[11pt,a4paper]{article}
\usepackage[utf8]{inputenc}
\usepackage[T1]{fontenc}
\usepackage[magyar]{babel}
\usepackage{amsmath,amssymb}
\usepackage{geometry}
\usepackage{tcolorbox}
\usepackage{xcolor}
\usepackage{tikz}

\geometry{left=25mm, right=25mm, top=25mm, bottom=25mm}

\definecolor{myprimary}{RGB}{41, 128, 185}
\definecolor{mybackground}{RGB}{245, 247, 250}
\definecolor{myboxborder}{RGB}{52, 152, 219}

\tcbset{
    colback=mybackground,
    colframe=myboxborder,
    boxrule=0.8pt,
    arc=4pt,
    left=10pt,
    right=10pt,
    top=10pt,
    bottom=10pt
}

\title{Kombináció}
\author{}
\date{}

%%% Megjelenítéshez automatikusan kiegészített csomagok %%%
\usepackage{amsfonts, amsthm}
\usepackage{pgfplots}
\pgfplotsset{compat=1.18}
\usepackage{graphicx}
\usepackage{enumitem}
\usepackage{booktabs}
\usepackage{float}
\usepackage{caption}
\usepackage{microtype}
\usepackage{abstract}

\usepackage{hyperref}

\begin{document}

\maketitle



\section{Mikor használjuk?}
Képzeld el, hogy a kedvenc pizzádat rendeled, és a 10-féle feltét közül ki kell választanod 3-at. Lényeges, hogy a szalámisat, a gombásat meg a kukoricásat milyen sorrendben mondod be a telefonba? Egyáltalán nem! Ugyanazt a pizzát kapod mindenesetre. 

Akkor nyúlunk a \textbf{kombináció} eszközéhez, amikor egy nagyobb halmazból kell kiválasztanunk egy kisebb csoportot úgy, hogy **a kiválasztás sorrendje NEM számít**, és elemeket **nem ismétlünk**. Ha a sorrend sem számít, és még ismételni is lehet (pl. gombás-gombás-kukorica), akkor az már egy másik történet (ismétléses kombináció), de először nézzük meg a klasszikus, alap esetet!

\section{Hogyan működik?}
Vegyünk egy nagyon egyszerű példát. Van $4$ darab különböző pizzafeltétünk: paprika ($P$), gomba ($G$), kukorica ($K$) és sajt ($S$). Ebből szeretnénk kiválasztani pontosan $2$-t a pizzánkra. Hányféleképpen tehetjük ezt meg?

Ha először a sorrendre is figyelnénk (ezt hívjuk \textbf{variációnak}), akkor azt mondanánk, hogy az első helyre választhatunk $4$-félét, a másodikra maradt $3$, azaz $4 \cdot 3 = 12$-féle sorrend jöhetne ki (pl. először a paprika, utána a gomba, az más, mint először a gomba, utána a paprika). 

De nekünk a pizzánkon a paprika-gomba és a gomba-paprika ugyanaz a feltétpár! Hányféleképpen tudunk sorba rendezni $2$ feltétet? Pontosan $2$-féleképpen ($2!$). Mivel minden egyes párost kétszer számoltunk meg (a két lehetséges sorrend miatt), ezért a lihegve kiszámolt $12$-es eredményt egyszerűen el kell osztanunk kettővel: $\frac{12}{2} = 6$.

Általánosan is ez a logika: először kiszámoljuk, hányféleképpen tudnánk sorba rendezni a kiválasztott elemeket, majd osztunk a sorrendek számával, hiszen a kombinációban a sorrend lényegtelen.

\section{Általános képlet és speciális esetek}
Ha $n$ darab különböző elemből szeretnénk kiválasztani $k$ darabet úgy, hogy a sorrend nem számít, akkor a következő \textbf{kombinációs képletet} használjuk (ezt úgy olvassuk ki, hogy "$n$ alatt a $k$"):

$$\binom{n}{k} = \frac{n!}{k!(n-k)!}$$

Ahol az $n!$ (n faktoriális) nem azt jelenti, hogy izgulunk, hanem azt, hogy $1 \cdot 2 \cdot 3 \dots \cdot n$.

\begin{tcolorbox}[title=Hasznos tulajdonságok és csapdák]
\begin{itemize}
    \item \textbf{Szimmetria szabály:} $\binom{n}{k} = \binom{n}{n-k}$. Ha 10 emberből kiválasztunk 8-at a csapatba, az ugyanannyi variáció, mintha azt mondanánk meg, hogy 2 embert hagyunk ki.
    \item \textbf{A klasszikus csapda:} Mindig tedd fel magadnak a kérdést: \textit{Ha felcserélem a kiválasztott elemek sorrendjét, akkor egy újabb egyedi esetet kaptam, vagy ugyanazt?} Ha ugyanazt, akkor kombináció! Ha számít a rang, a helyezés vagy a sorrend, akkor ez bizony nem kombináció.
\end{itemize}
\end{tcolorbox}

\section{Dolgozat- / Vizsga-gyorsteszt (Döntési fa)}
Amikor olvasod a feladatot a dolgozatban, futtasd le ezt a 3 kérdést a fejedben:
\begin{enumerate}
    \item \textbf{Minden elemet felhasználunk?} Ha csak egy részét válogatjuk ki (pl. 10-ből 3-at), akkor ez nem permutáció.
    \item \textbf{Számít a sorrend?} (Pl. elnök és alelnök választás vs. egyszerű ötfős bizottság). Ha \textbf{nem} számít a sorrend, akkor jó helyen jársz, ez kombináció!
    \item \textbf{Visszatehetjük-e az elemet?} (Ismétlés). Ha a szöveg nem említi az ismétlést, és a tárgyakból/emberekből csak egy van, akkor sima (ismétlés nélküli) kombinációról van szó.
\end{enumerate}

\section{Tipikus és trükkös feladatok levezetése}

\subsection*{Feladat: A klasszikus bizottság-választás}
\begin{tcolorbox}
Egy osztályban $12$ fiú és $8$ lány tanul. Hányféleképpen lehet 5 fős bizottságot küldeni a diákönkormányzatba, ha a bizottságban \textbf{pontosan 3 fiúnak} kell lennie?
\end{tcolorbox}

\textbf{Gondolatmenet és lépések:}
1. \textbf{Mit látunk?} Összesen $20$ tanulónkból kell kiválasztani $5$-öt, de van egy szigorú megkötés: a csapatnak pontosan $3$ fiúból és így ugyebár $5 - 3 = 2$ lányból kell állnia.
2. \textbf{Miért így gondolkodunk?} Külön kell kiválasztanunk a fiúkat a fiúk csoportjából, és a lányokat a lányok csoportjából. Mivel a bizottságon belül a sorrend mindegy (tök mindegy, hogy Jancsi vagy Pistike lép be először a terembe), ez kombináció.
3. \textbf{A fiúk kiválasztása:} $12$ fiúból kell $3$-at kiválasztani. Ennek a lehetőségei:
$$\binom{12}{3} = \frac{12 \cdot 11 \cdot 10}{1 \cdot 2 \cdot 3} = 220$$
4. \textbf{A lányok kiválasztása:} $8$ lányból kell $2$-t kiválasztani. Ennek a lehetőségei:
$$\binom{8}{2} = \frac{8 \cdot 7}{1 \cdot 2} = 28$$
5. \textbf{Összeadás vagy szorzás?} Ez egy klasszikus hibaforrás! Sokan összeadják a lehetőségeket. De figyelj: minden egyes fiúcsapathoz \textit{bármelyik} lánycsapat tartozhat. Ezért a kombinatorikai alapelv (szorzási elv) szerint a két eredményt \textbf{össze kell szorozni}:
$$\text{Összesen} = 220 \cdot 28 = 6160$$
\textbf{Megoldás:} A bizottságot $6160$-féleképpen lehet összeállítani.

\subsection*{Feladat: Geometriai kombináció}
\begin{tcolorbox}
Egy konvex tízszög csúcsai közül kiválasztunk hármat. Hány olyan háromszög van, amelynek a csúcsai a tízszög adott csúcsai közül kerülnek ki?
\end{tcolorbox}

\textbf{Gondolatmenet és lépések:}
1. \textbf{Mit látunk?} Van $10$ darab pontunk (a tízszög csúcsai), és ezek közül kell $3$-at összekötnünk, hogy háromszöget kapjunk.
2. \textbf{Hol a csapda?} A diákok gyakran megijednek, mert ez geometriai feladatnak tűnik, pedig a lényege tiszta kombináció. Ha kiválasztunk pl. az $A$, $B$ és $C$ csúcsokat, azzal máris meghatároztunk egy háromszöget. Számít-e, hogy $ABC$, $BCA$ vagy $CAB$ sorrendben nevezzük meg a csúcsokat? Nem, a háromszög ugyanaz marad!
3. \textbf{A számítás:} $10$ pontból kell kiválasztanunk $3$-at, sorrend nélkül:
$$\binom{10}{3} = \frac{10 \cdot 9 \cdot 8}{1 \cdot 2 \cdot 3} = 120$$
\textbf{Megoldás:} Pontosan $120$ ilyen háromszög létezik. (Nincs is annyi sok, ugye?)

\end{document}